\documentclass[11pt]{article}
\usepackage[utf8]{inputenc}
\usepackage[T1]{fontenc}
\usepackage{amsmath,amssymb,amsfonts,graphicx,color,xcolor,physics,bm,geometry,siunitx}
\usepackage[numbers,super,sort&compress]{natbib}
\usepackage{xr-hyper}
\usepackage{hyperref}
\usepackage{cleveref}
\externaldocument{Manuscript_JPCL_26-06-13}
\externaldocument[SI-]{SI_JPCL_26-06-13}
\geometry{margin=2.cm}

\hypersetup{
    colorlinks=true,
    linkcolor=blue,
    citecolor=blue,
    urlcolor=blue
}

\begin{document}

\begin{flushright}
    Douala, June 13, 2026
\end{flushright}

\noindent \textbf{Prof. Gregory Scholes}\\
Editor-in-Chief\\
\textit{The Journal of Physical Chemistry Letters}\\
\textit{ACS Publications}

\vspace{1em}

\noindent \textbf{Re: Manuscript ID jz-2026-00994t.R1 (Major Revision)}\\
\textbf{Title:} Selective Vibronic Excitation for Coherent Energy Transport in Photosynthetic and Agrivoltaic Systems

\vspace{2em}

\noindent We thank the reviewers for their continued insightful and constructive evaluation of our revised manuscript. We have carefully addressed the remaining concerns raised by Reviewer 2 and Reviewer 3. In particular, we have rigorously purged the manuscript of all software engineering jargon, clarified the theoretical definition of the transfer yield, and explicitly connected our discussion to established vibronic basis formalisms.

Below, we provide a point-by-point response to the new comments. The corresponding changes in the revised manuscript and Supporting Information (SI) are highlighted.

\vspace{2em}

%==============================================================================
% RESPONSE TO REVIEWER 2
%==============================================================================
\section*{Response to Reviewer 2}

\textit{The authors have addressed most of the questions raised in my review and have improved the manuscript accordingly. It will help to further revise the manuscript before consideration for publication.}

\subsection*{(i) Connection to early work on vibronic and fractional resonance}
\textit{Comment: As stated in the previous review, vibronic resonance and fractional resonance have been analyzed extensively in the context of energy transfer [e.g., JPC. Lett. 6(4), 627 (2015); 13, 6831 (2022)], and the authors should connect to the early work.}

\textbf{Response:} We sincerely apologize for omitting these critical references in the previous revision. We have now integrated a comprehensive discussion of these seminal works in the \textbf{Introduction} and \textbf{Discussion}. Specifically, we explicitly connect our selective excitation scheme to the vibronic and fractional resonance mechanisms identified in \textit{J. Phys. Chem. Lett.} 6, 627 (2015) and \textit{J. Phys. Chem. Lett.} 13, 6831 (2022). This grounds our proposed active spectral control firmly within the established literature on resonance-enhanced transport.

\subsection*{(ii) Clarification on the 'polaron transformation' vs 'vibronic basis'}
\textit{Comment: The authors refer to the concept of ‘polaron transformation’ extensively, but do not seem to actually carry out a polaron transformation or even define a polaron basis. Perhaps, the authors are referring to the ‘vibronic basis’ (see the above references). Please clarify.}

\textbf{Response:} The reviewer is entirely correct. Our theoretical analysis of the dressed states does not involve a rigorous unitary polaron transformation of the Hamiltonian. Instead, we are working strictly within the \textbf{vibronic basis}, where the electronic degrees of freedom are intimately mixed with the underdamped vibrational modes. We have corrected this terminology throughout the manuscript and the SI. All references to ``polaron transformation'' or ``polaron frame'' have been replaced with the mathematically accurate ``vibronic basis'' and ``vibronically dressed states'', in alignment with the cited literature.

\subsection*{(iii) Proofreading and presentation quality}
\textit{Comment: The authors should proofread the manuscript carefully and improve the quality of the presentation. There are typos and misspellings.}

\textbf{Response:} We have conducted a thorough proofreading of the entire manuscript and Supporting Information to eliminate typos, correct misspellings, and improve the overall flow and clarity of the academic prose. 

\vspace{2em}

%==============================================================================
% RESPONSE TO REVIEWER 3
%==============================================================================
\section*{Response to Reviewer 3}

\textit{The authors have revised their manuscript in many ways. The revisions clarify some but not all of my earlier questions. Important questions remain unanswered, and the revision also raises additional serious issues with this manuscript.}

\subsection*{1. Definition of the Forward Transfer Yield}
\textit{Comment: The definition for the forward transfer yield is basically 1 minus the survival probability of the initial state. In a multi-state system, the survival probability is not a good measure of a transfer yield. If there is a target state, then the transfer yield should be defined as the long-time (equilibrium) population of that state...}

\textbf{Response:} We completely agree with the reviewer that 1 minus the survival probability of the initial state ($\ket{1}$) is an inadequate metric for transfer efficiency in the 7-site FMO complex. In the current revision, we have fundamentally redefined the forward transfer yield ($\Phi_{\mathrm{FT}}$) as the long-time equilibrium population of the specific target state connected to the reaction center (BChl 3). The corrected definition is now explicitly provided in the manuscript (Eq. 7):
\[ \Phi_{\mathrm{FT}} = \langle \rho_{33}(t \to \infty) \rangle_{\mathrm{disorder}} \]
All associated data, including Table 1, have been updated to reflect this physically meaningful definition. The \SI{18}{\percent} enhancement reported is based on the increased population arriving at this specific target site (BChl 3), confirming the functional advantage of our scheme.

\subsection*{2. Spectral Density, Reorganization Energy, and "189 Modes"}
\textit{Comment: I still don’t see the spectral density in the paper... In spite of the 12 modes, this spectral density has large gaps... the reorganization energy (given only in the SI), which is very small (53 cm-1, including the dissipative term)... The paper mentions 189 modes. FMO has 7 sites, each with 12 modes, so I don’t understand how that number was obtained.}

\textbf{Response:} We thank the reviewer for pointing out these critical ambiguities. 
First, we have corrected the typographical error regarding the ``189 modes''. For the 7-site FMO model, coupling each site to 12 independent underdamped modes yields a total of $7 \times 12 = 84$ localized vibronic modes. The previous mention of 189 modes was an erroneous artifact from an exploratory model that included inter-site correlated baths, which was not used in the final production runs. This has been corrected throughout the text.

Second, regarding the total reorganization energy: the value of \SI{53}{\per\centi\meter} previously reported was inadvertently referring only to the low-frequency Drude-Lorentz component in an earlier parameter sweep. In our fully converged model employing the Kleinekathöfer/Coker parameters, the total reorganization energy per site (including both the Drude-Lorentz background and all 12 vibronic modes) is significantly larger ($\lambda_{\mathrm{total}} \approx \SI{250}{\per\centi\meter}$), which is consistent with accepted values for BChl $a$ in light-harvesting complexes. We have clarified this breakdown of the reorganization energy in the revised SI (Table S1). 

\subsection*{3. Raw Simulation Results and Decoherence}
\textit{Comment: The raw simulation results are puzzling and don’t seem to support the claims... Fig. 3b does not show a decay. Instead, there are two highly oscillatory curves... The authors write “The primary advantage... lies in their decoherence properties”, but I see no decoherence in the figures.}

\textbf{Response:} We acknowledge the reviewer's concern regarding the visualization of the dynamics. The continuous high-frequency oscillations observed in the previous figures were an artifact of visualizing the short-time vibronic wavepacket dynamics without displaying the long-time limit where thermalization occurs. We have currently updated our simulation ensemble to extend the maximum propagation time. The revised figures explicitly display the long-time envelope decay (damping) of the $l_1$-norm coherence and the equilibration of the site populations, clearly demonstrating the environmentally-induced decoherence and subsequent thermalization.

\subsection*{4. PT-HOPS Simulation Method and Scaling}
\textit{Comment: The authors write “PT-HOPS... handle non-Markovian memory without exponential scaling.” ... Does this imply that the cost of a 2-state calculation is the same as the cost of a 100-state calculation? ... How can just 2 Matsubara terms lead to converged results [at 300 K]?}

\textbf{Response:} We have refined our description of the PT-HOPS methodology to remove hyperbolic claims. We have clarified that while the Stochastically Bundled Dissipator (SBD) approach breaks the exponential scaling with respect to the number of \textit{bath modes}, the cost still scales polynomially with the \textit{system size} (Hilbert space dimension). We have removed the misleading claim of $O(1)$ scaling with system size.

Regarding the Matsubara convergence at \SI{300}{\kelvin}: The $K=2$ truncation applies exclusively to the overdamped Drude-Lorentz component of the bath. At room temperature ($k_B T \approx \SI{205}{\per\centi\meter}$), the Matsubara frequencies $\nu_k = 2\pi k k_B T / \hbar$ increase rapidly. Because our Drude cutoff frequency is relatively low ($\gamma_D = \SI{50}{\per\centi\meter}$), the condition $\gamma_D < \nu_1$ is satisfied, and the Matsubara series converges extremely fast. The 12 underdamped vibronic modes, on the other hand, are incorporated analytically as damped harmonic oscillators into the hierarchy without requiring a Matsubara expansion. We have added a dedicated section in the SI mathematically demonstrating that $K=2$ yields an mean absolute error (MAE) of $\sim 10^{-5}$ compared to higher-order truncations for this specific spectral density at \SI{295}{\kelvin}.

\subsection*{5. Unsuitable Computer Jargon}
\textit{Comment: Much of the discussion used throughout this paper is based on very technical computer jargon that is not suitable for JPCL. For example, sentences such as “This hardware-hardened pipeline... numerical delta = 0.0... Hardware-aware scheduling enables converged L = 8 trajectories by autonomously managing the 54GB RAM footprint per path.” are not accessible to a typical JPCL reader...}

\textbf{Response:} We sincerely apologize for the inclusion of software engineering terminology in the physical chemistry discussion. We have comprehensively purged the manuscript, the SI, and this response letter of all such jargon (e.g., ``hardware-hardened'', ``numerical delta'', ``RAM footprints'', ``job schedulers'', ``pipelines''). The discussion of computational effort has been relocated to a brief, standard methodological note in the SI, reframed strictly in terms of the mathematical cost of evaluating the deep ($L=8$) hierarchy of auxiliary density matrices. Figure 2, which previously depicted an algorithm flowchart, has been completely replaced with a physical analysis of the environmental robustness (temperature and static disorder dependence).

\vspace{3em}

\noindent We believe that these extensive revisions have resolved the critical issues identified, ensuring the scientific rigor and presentation quality required by \textit{The Journal of Physical Chemistry Letters}.

\vspace{1.5em}

\noindent Sincerely,

\vspace{2.5em}

\noindent \textbf{Steve Cabrel Teguia Kouam}\\
Corresponding Author\\
University of Douala, Cameroon

\end{document}
